\section{Exact arithmetic certificates for numerical examples}
\label{sec:certificates}

The numerical examples are substitutions into the proved mean and variance
formulas. The exact fractions below, together with the rounding lemma,
provide arithmetic proofs of the displayed decimals. They do not depend on
floating-point output, simulation, or the software-verification results.

\subsection{Exact values behind the IQ-150 examples}

\begin{corollary}[Exact IQ means and variances]
\label{cor:exact-iq-values}
Under the joint Gaussian model, let \(r,k\geq1\) be integers and set
\(c=3/5\), \(a=10/3\), and \(N=r+k\). On the IQ reporting scale
\(Q=100+15X\), write \(V_{\mathrm{pop}}\) and \(V_{\mathrm{cond}}\) for the
population variance and the variance conditional on all \(r\) initial
IQ scores being \(150\), respectively, of the actual IQ-score mean. Then
\begin{align}
 m_{r,k}^{\mathrm{IQ}}
   &=100+\frac{50r(3N+2)}{N(3r+2)},
      \label{eq:exact-iq-mean}\\
 V_{\mathrm{pop}}
   &=\frac{45(3N+2)}{N},
      \label{eq:exact-iq-population-variance}\\
 V_{\mathrm{cond}}
   &=\frac{90k(3N+2)}{N^2(3r+2)}.
      \label{eq:exact-iq-conditional-variance}
\end{align}
\end{corollary}

\begin{proof}
The definitions give
\[
 d_r=1+(r-1)\frac35=\frac{3r+2}{5},\qquad
 d_N=1+(N-1)\frac35=\frac{3N+2}{5}.
\]
Substitution into \eqref{eq:iq150-mean} gives
\(100+50r d_N/(N d_r)\), which is
\eqref{eq:exact-iq-mean}. Squaring
\eqref{eq:iq150-population-sd} gives \(225d_N/N\), hence
\eqref{eq:exact-iq-population-variance}. Squaring
\eqref{eq:iq150-conditional-sd} gives
\[
 \frac{225k(1-3/5)d_N}{N^2d_r}
   =\frac{90k(3N+2)}{N^2(3r+2)},
\]
which is \eqref{eq:exact-iq-conditional-variance}. The affine moment
identities in Theorem~\ref{thm:moments} justify the factor \(15^2=225\)
for variances and the additive \(100\) for means.
\end{proof}

For the nine parameter pairs in Section~\ref{sec:iq150-examples}, the
resulting exact quantities are as follows. The SDs are the nonnegative
square roots of the two variance columns.

\begin{center}
\begin{tabular}{rrr r r r}
\hline
\(r\) & \(k\) & \(N\) & Exact IQ mean & \(V_{\mathrm{pop}}\) &
\(V_{\mathrm{cond}}\) \\
\hline
1 & 1  & 2  & \(140\) & \(180\) & \(36\) \\
1 & 2  & 3  & \(410/3\) & \(165\) & \(44\) \\
1 & 10 & 11 & \(1450/11\) & \(1575/11\) & \(6300/121\) \\
\hline
2 & 1  & 3  & \(875/6\) & \(165\) & \(55/4\) \\
2 & 2  & 4  & \(575/4\) & \(315/2\) & \(315/16\) \\
2 & 10 & 12 & \(1675/12\) & \(285/2\) & \(475/16\) \\
\hline
3 & 1  & 4  & \(1625/11\) & \(315/2\) & \(315/44\) \\
3 & 2  & 5  & \(1610/11\) & \(153\) & \(612/55\) \\
3 & 10 & 13 & \(20450/143\) & \(1845/13\) & \(36900/1859\) \\
\hline
\end{tabular}
\end{center}

For example, the last row has \(r=3\), \(k=10\), and \(N=13\). Direct
integer substitution yields
\[
 100+\frac{50\cdot3\cdot41}{13\cdot11}
    =\frac{20450}{143},\qquad
 \frac{45\cdot41}{13}=\frac{1845}{13},\qquad
 \frac{90\cdot10\cdot41}{13^2\cdot11}
    =\frac{36900}{1859}.
\]
The other rows follow by the same substitutions; all quantities in this
table are exact rational numbers, rather than decimal approximations.

\subsection{A proof of the decimal rounding rule}

\begin{lemma}[Exact rounding test]
\label{lem:rounding}
Let \(T=10^p\), where \(p\) is a positive integer, and let \(t\) be an
integer. Rounding a real number to \(p\) decimal places means choosing a
nearest point of the grid \(\{j/T:j\in\mathbb Z\}\).
\begin{enumerate}
\item For \(x=P/Q\), with integers \(P,Q\) and \(Q>0\), define
\begin{equation}
 D_T=Q-2|TP-tQ|.
 \label{eq:rounding-mean-certificate}
\end{equation}
If \(D_T>0\), the unique nearest grid point to \(x\) is \(t/T\).
\item For a variance \(v=P/Q\geq0\), with integers \(P,Q\) and \(Q>0\),
and an integer \(t\geq1\), define
\begin{align}
 L_T&=4T^2P-(2t-1)^2Q,\notag\\
 U_T&=(2t+1)^2Q-4T^2P.
 \label{eq:rounding-sd-certificates}
\end{align}
If \(L_T>0\) and \(U_T>0\), the unique nearest grid point to
\(\sqrt v\) is \(t/T\).
\end{enumerate}
All these strict inequalities exclude a halfway tie.
\end{lemma}

\begin{proof}
For the first part, \(D_T>0\) is equivalent to
\[
 |P/Q-t/T|=\frac{|TP-tQ|}{TQ}<\frac1{2T}.
\]
Every other grid point is at distance at least \(1/T\) from \(t/T\), so
the triangle inequality makes its distance from \(x\) strictly greater
than \(1/(2T)\). Thus \(t/T\) is the unique nearest grid point.

For the second part, divide the two positive integer differences by
\(4T^2Q\). They give
\[
 \left(\frac{2t-1}{2T}\right)^2
     <v<
 \left(\frac{2t+1}{2T}\right)^2.
\]
Both bounding numbers before squaring are positive. Squaring is strictly
increasing on the nonnegative real numbers, so these inequalities imply
\[
 \frac{t-1/2}{T}<\sqrt v<\frac{t+1/2}{T}.
\]
Therefore \(|\sqrt v-t/T|<1/(2T)\), and the grid-distance argument in
the first part proves the result. This uses only the order properties of
the real numbers and the definition of nearest-decimal rounding.
\end{proof}

\subsection{Positive integer certificates for every IQ table row}

For each row of the displayed IQ table, take \(T=100\). Use the reduced
fraction \(P/Q\) in the exact-value table above. For the mean, \(t\) is
100 times its displayed two-decimal value. For each SD, \(t\) is 100 times
its displayed value, while \(P/Q\) is its exact \emph{variance}.
The following table gives the integer differences from
\eqref{eq:rounding-mean-certificate} and
\eqref{eq:rounding-sd-certificates}. In the SD formulas,
\(4T^2=40000\).

\begin{center}
\begingroup
\small
\setlength{\tabcolsep}{5pt}
\begin{tabular}{rr r rr rr}
\hline
\(r\) & \(k\) & Mean \(D_{100}\) &
\(L_{\mathrm{pop}}\) & \(U_{\mathrm{pop}}\) &
\(L_{\mathrm{cond}}\) & \(U_{\mathrm{cond}}\) \\
\hline
1 & 1  & 1  & 1511   & 9225  & 2399   & 2401 \\
1 & 2  & 1  & 239    & 10041 & 4375   & 929 \\
1 & 10 & 7  & 9061   & 96275 & 47871  & 651025 \\
2 & 1  & 2  & 239    & 10041 & 3676   & 8196 \\
2 & 2  & 4  & 9838   & 10242 & 11696  & 45136 \\
2 & 10 & 4  & 4462   & 14642 & 25264  & 44496 \\
3 & 1  & 5  & 9838   & 10242 & 6100   & 88236 \\
3 & 2  & 3  & 4271   & 5625  & 11105  & 135855 \\
3 & 10 & 57 & 100907 & 22957 & 175221 & 6457691 \\
\hline
\end{tabular}
\endgroup
\end{center}

Every listed integer is strictly positive. Lemma~\ref{lem:rounding}
therefore proves each of the nine mean roundings and all eighteen SD
roundings in the IQ table. For example, in the \(r=3,k=10\) row the
displayed conditional SD is \(4.46\), so \(t=446\) and \(v=36900/1859\).
The exact differences are
\[
 40000\cdot36900-891^2\cdot1859=175221>0,
\]
\[
 893^2\cdot1859-40000\cdot36900=6457691>0.
\]
Thus the true SD is strictly between \(4.455\) and \(4.465\), proving
that its nearest two-decimal representation is \(4.46\).

\subsection{Certificates for the other numerical examples}

For the earlier standardized examples, substitute \(a=2\), \(c=3/5\),
and \(k=3\) into
\eqref{eq:multiple-predictions}, \eqref{eq:multiple-mean},
\eqref{eq:multiple-population-sd}, and
\eqref{eq:multiple-conditional-sd}. With \(N=r+3\), the exact quantities
are
\begin{center}
\begin{tabular}{r r r r r}
\hline
\(r\) & Each prediction \(b_r\) & Inclusive mean \(B_{r,3}\) &
Population variance & Conditional variance \\
\hline
1 & \(6/5\) & \(7/5\) & \(7/10\) & \(21/100\) \\
2 & \(3/2\) & \(17/10\) & \(17/25\) & \(51/500\) \\
3 & \(18/11\) & \(20/11\) & \(2/3\) & \(2/33\) \\
\hline
\end{tabular}
\end{center}
In particular, the terminating values \(1.2\), \(1.4\), \(1.5\), and
\(1.7\), and their zero-padded forms, equal the corresponding fractions
exactly. Each zero SD stated for a fixed recursive mean follows from
Theorem~\ref{thm:moments}, rather than a rounding approximation.

The nonterminating means and predictions used elsewhere have the following
certificates. Here \(P/Q\) denotes the exact value itself, and each positive
\(D_T\) is computed by \eqref{eq:rounding-mean-certificate}.

\begin{center}
\begin{tabular}{l r r r r}
\hline
Quantity & \(P/Q\) & \(T\) & Displayed value & \(D_T\) \\
\hline
\(b_3\), standardized example & \(18/11\) & \(10^6\) & \(1.636364\) & 3 \\
\(B_{3,3}\), standardized example & \(20/11\) & \(10^6\) & \(1.818182\) & 7 \\
\(q_3\), IQ-150 prediction & \(1550/11\) & \(10^6\) & \(140.909091\) & 9 \\
\(q_3\), IQ-150 prediction & \(1550/11\) & \(100\) & \(140.91\) & 9 \\
Worked IQ mean, \(r=2,k=10\) & \(1675/12\) & \(10^6\) & \(139.583333\) & 4 \\
\hline
\end{tabular}
\end{center}
The other IQ predictions, \(q_1=130\) and \(q_2=275/2=137.5\), are exact
terminating values.

Finally, the following positive integer differences prove all displayed
six-decimal SD values. In this table \(T=10^6\), \(P/Q\) is the exact
variance, and \(L_T,U_T\) are given by
\eqref{eq:rounding-sd-certificates}.

\begin{center}
\begingroup
\small
\setlength{\tabcolsep}{4pt}
\begin{tabular}{l r r r r}
\hline
SD quantity & Variance \(P/Q\) & Displayed SD & \(L_T\) & \(U_T\) \\
\hline
\(r=1\), population & \(7/10\) & \(0.836660\) & 35242390 & 31690410 \\
\(r=1\), conditional & \(21/100\) & \(0.458258\) & 25477500 & 341128900 \\
\(r=2\), population & \(17/25\) & \(0.824621\) & 103097975 & 61826225 \\
\(r=2\), conditional & \(51/500\) & \(0.319374\) & 1134995500 & 142500500 \\
\(r=3\), population & \(2/3\) & \(0.816497\) & 1585853 & 18010075 \\
\(r=3\), conditional & \(2/33\) & \(0.246183\) & 31323575 & 33668737 \\
Worked IQ, population & \(285/2\) & \(11.937336\) & 169283518 & 21713858 \\
Worked IQ, conditional & \(475/16\) & \(5.448624\) & 125135856 & 572288016 \\
\hline
\end{tabular}
\endgroup
\end{center}
Every entry in the last two certificate columns is strictly positive, so
Lemma~\ref{lem:rounding} proves the claimed six-decimal roundings. The exact
fractions also establish the worked IQ-row variances \(285/2\) and
\(475/16\). All decimal claims in the numerical examples are therefore
either exact terminating fractions or consequences of these strict rational
inequalities.
